% mosaic_3D_Plane.tex
% Preamble (main .tex file):
%   \usepackage{pgfplots}
%   \pgfplotsset{compat=1.18}
%   \usepgfplotslibrary{colormaps}

\begin{tikzpicture}
  % --------- global resolution knob ----------
  % Change this single integer to tune resolution everywhere
  \pgfmathtruncatemacro{\NsampBase}{16}   % e.g. 8 (coarse), 16 (medium), 32 (fine)

  % Derived resolutions
  \pgfmathtruncatemacro{\Nsamp}{\NsampBase}          % angular resolution on spheres
  \pgfmathtruncatemacro{\NsampY}{0.75*\NsampBase}    % polar samples on spheres
  \pgfmathtruncatemacro{\NsampArc}{7.5*\NsampBase}   % samples along intersection arcs

  % Mosaic kernel parameters (same as pv_arc)
  \pgfmathsetmacro{\sigmaG}{0.2964692365} % Gaussian sigma [rad]
  \pgfmathsetmacro{\gammaL}{0.3490658504} % Lorentzian HWHM [rad]
  \pgfmathsetmacro{\eta}{0.3}             % Lorentzian weight in ω
  \pgfmathsetmacro{\thetaZero}{0}        % Bragg-circle polar angle θ0 (deg)

  % Sphere geometry
  \pgfmathsetmacro{\R}{1.0}               % Bragg-sphere radius
  \pgfmathsetmacro{\Rext}{1.3}            % Axis length
  \pgfmathsetmacro{\eps}{0.03}            % Small z-offset for G_y arrow
  \pgfmathsetmacro{\Rdiag}{\R/sqrt(2)}    % |(+Gz - Gy)| projected on (x,z)

  % Intersection circle of two equal spheres (both radius R):
  % Centers S1 = (0,0,0) and S2 = (-R/√2, 0, R/√2)
  % Circle center C and radius r_circ:
  \pgfmathsetmacro{\cx}{-0.5*\R/sqrt(2)}      % circle center x  (Gy)
  \pgfmathsetmacro{\cz}{ 0.5*\R/sqrt(2)}      % circle center z  (Gz)
  \pgfmathsetmacro{\rcirc}{0.5*sqrt(3)*\R}    % circle radius

  % One-color colormap for auxiliary sphere (grey)
  \pgfplotsset{
    colormap={auxGrayMap}{
      rgb255(0pt)   =(160,160,160);
      rgb255(1000pt)=(160,160,160)
    }
  }

  \begin{axis}[
    view={-35}{20},
    axis equal image,
    hide axis,
    shader=interp,
    colormap/viridis,
    point meta rel=per plot,
    title={Surface density $\rho(\theta)$ on the Bragg sphere},
    title style={yshift=8pt},
    unbounded coords=jump,     % break curves when points are out of range
  ]

    % ---------------- MAIN BRAGG SPHERE, BACK HALF (x + y > 0) ----------------
    \addplot3[
      surf,
      domain=0:360,      % azimuth φ
      y domain=0:90,     % polar angle θ (from +z)
      samples=\Nsamp,
      samples y=\NsampY,
      variable=\u,
      variable y=\v,
      opacity=0.8,
      draw=none,
      point meta={
        (
          \eta * (1/pi * \gammaL /
                  (((mod(\v-\thetaZero+180,360)-180)*pi/180)^2 + \gammaL^2))
        )
        +
        (
          (1-\eta) * (1/(sqrt(2*pi)*\sigmaG) *
            exp(-0.5*(((mod(\v-\thetaZero+180,360)-180)*pi/180)/\sigmaG)^2))
        )
      },
      restrict expr to domain={x + y}{0:10},  % x + y > 0  -> back half
    ]
    ({cos(\u)*sin(\v)}, {sin(\u)*sin(\v)}, {cos(\v)});

    % ---------------- FULL AUXILIARY SPHERE (grey) ----------------
    % Center: (-Rdiag, 0, Rdiag), radius = R
    \addplot3[
      surf,
      domain=0:360,
      y domain=0:180,    % full sphere
      samples=\Nsamp,
      samples y=\NsampY,
      variable=\U,
      variable y=\V,
      shader=interp,
      draw=none,
      opacity=0.6,
      point meta=0,
      colormap name=auxGrayMap,
    ]
    ({- \Rdiag + \R * cos(\U)*sin(\V)},
     {0        + \R * sin(\U)*sin(\V)},
     {  \Rdiag + \R * cos(\V)});

    % ---------------- MAIN BRAGG SPHERE, FRONT HALF (x + y < 0) ----------------
    \addplot3[
      surf,
      domain=0:360,      % azimuth φ
      y domain=0:90,     % polar angle θ (from +z)
      samples=\Nsamp,
      samples y=\NsampY,
      variable=\u,
      variable y=\v,
      opacity=0.8,
      draw=none,
      point meta={
        (
          \eta * (1/pi * \gammaL /
                  (((mod(\v-\thetaZero+180,360)-180)*pi/180)^2 + \gammaL^2))
        )
        +
        (
          (1-\eta) * (1/(sqrt(2*pi)*\sigmaG) *
            exp(-0.5*(((mod(\v-\thetaZero+180,360)-180)*pi/180)/\sigmaG)^2))
        )
      },
      restrict expr to domain={x + y}{-10:0}, % x + y < 0  -> front half
    ]
    ({cos(\u)*sin(\v)}, {sin(\u)*sin(\v)}, {cos(\v)});

    % ---------------- RECIPROCAL-SPACE AXES ----------------
    % G_z
    \addplot3[->,thick] coordinates {(0,0,0) (0,0,\Rext)};
    \node[anchor=south] at (axis cs:0,0,\Rext+0.05) {$G_z$};

    % G_x: along negative y (dashed inside, solid outside)
    \addplot3[thick,dashed] coordinates {(0,0,0) (0,-\R,0)};
    \addplot3[->,thick]      coordinates {(0,-\R,0) (0,-\Rext,0)};
    \node[anchor=north] at (axis cs:0,-\Rext-0.05,0) {$G_x$};

    % G_y: along positive x (dashed inside, solid outside, slightly lifted in z)
    \addplot3[thick,dashed] coordinates {(0,0,\eps) (\R,0,\eps)};
    \addplot3[<-,thick]      coordinates {(\R,0,\eps) (\Rext,0,\eps)};
    \node[anchor=west] at (axis cs:\Rext+0.05,0,\eps) {$G_y$};

    % +G_z - G_y vector
    \addplot3[<-,thick,oiOrange]
      coordinates {(0,0,0) (-\Rdiag,0,\Rdiag)};
    \node[anchor=south east]
      at (axis cs:-\Rdiag-0.05,0,\Rdiag+0.05)
      {$+G_z - G_y$};

    % ---------------- INTERSECTION CIRCLE PARAMETERISATION ----------------
    % x(t) = cx + rcirc*(1/√2)*sin t   (Gy)
    % y(t) = rcirc*cos t               (−Gx)
    % z(t) = cz + rcirc*(1/√2)*sin t   (Gz)
    % So:  G_x > 0  ⇔  y < 0,   G_x < 0 ⇔ y > 0.

    % BACK ARC: G_x < 0 (y > 0), G_z >= 0
    \addplot3[
      thick,
      color=red,
      opacity=0.6,
      samples=\NsampArc,
      domain=0:360,
      variable=\t,
      restrict expr to domain={z}{0:10},    % G_z >= 0
      restrict expr to domain={y}{0:10},    % G_x < 0  (y > 0)
    ]
    ({
       \cx + \rcirc*( 1/sqrt(2) * sin(\t) )
     },
     {
       0   + \rcirc*(            cos(\t) )
     },
     {
       \cz + \rcirc*( 1/sqrt(2) * sin(\t) )
     });

    % FRONT ARC: G_x > 0 (y < 0), G_z >= 0
    \addplot3[
      very thick,
      color=red,
      samples=\NsampArc,
      domain=0:360,
      variable=\t,
      restrict expr to domain={z}{0:10},     % G_z >= 0
      restrict expr to domain={y}{-10:0},    % G_x > 0  (y < 0)
    ]
    ({
       \cx + \rcirc*( 1/sqrt(2) * sin(\t) )
     },
     {
       0   + \rcirc*(            cos(\t) )
     },
     {
       \cz + \rcirc*( 1/sqrt(2) * sin(\t) )
     });

  \end{axis}
\end{tikzpicture}
